\documentclass[11pt,a4paper]{article}
\usepackage[utf8]{inputenc}
\usepackage[portuguese]{babel}
\usepackage{amsmath}
\usepackage{amssymb}
\usepackage{graphicx}
\usepackage{listings}
\usepackage{xcolor}
\usepackage{float}
\usepackage{geometry}
\usepackage{fancyhdr}
\usepackage{hyperref}
\usepackage{array}
\usepackage{booktabs}

\geometry{left=1.5cm,right=1.5cm,top=1.5cm,bottom=1.5cm}
\setlength{\parskip}{0.3em}
\pagestyle{fancy}

\lhead{CC8550 --- Simulação e Teste de Software}
\chead{Aula 05}
\rhead{2º Semestre 2025}
\cfoot{Página \thepage}

\title{
    {\Large \textbf{Simulação e Teste de Software}} \\
    {\Large \textbf{Aula 05 --- Teste de Mutação (Caixa Eletrônico)}}
}

\author{
    \textbf{Aluno:} Lucas Rebouças Silva \\
    \textbf{R.A.:} 22.122.048-6 \\[0.3cm]
    \textbf{Professor:} Luciano Rossi \\
    \textbf{Centro Universitário FEI} \\
    Ciência da Computação
}

\date{2º Semestre de 2025}

\begin{document}

\maketitle

\newpage
\tableofcontents
\newpage

\section{Introdução}

Neste laboratório foi desenvolvido um programa simples de \emph{caixa eletrônico}
com cinco funcionalidades principais e um conjunto de testes automatizados.
Foram aplicadas técnicas de teste de mutação manual e automática (\texttt{mutmut})
para avaliar a eficácia da suíte de testes.

\section{Descrição do Programa}

O sistema de caixa eletrônico foi implementado no arquivo
\texttt{caixa\_eletronico.py}. As cinco funcionalidades principais são:

\begin{enumerate}
  \item Consulta de saldo (\texttt{consultar\_saldo});
  \item Depósito em conta (\texttt{depositar});
  \item Saque em conta (\texttt{sacar});
  \item Transferência entre contas (\texttt{transferir});
  \item Emissão de extrato (\texttt{extrato}).
\end{enumerate}

\section{Estratégia de Testes}

Os testes foram escritos com \texttt{pytest} no arquivo
\texttt{test\_caixa\_eletronico.py}. Foram especificados dez casos de teste,
cobrindo depósitos válidos e inválidos, saques válidos e inválidos,
transferências e extrato.

\section{Mutantes Manuais}

\subsection*{O que são mutantes mortos e sobreviventes?}

Um \textbf{mutante morto} é uma versão do programa em que foi introduzida uma alteração (mutação) e, ao rodar a suíte de testes, \emph{pelo menos um teste falha}. Isso indica que os testes detectaram a diferença entre o código original e o mutado: o defeito simulado foi revelado e o mutante é considerado ``morto'' pela suíte. Já um \textbf{mutante sobrevivente} é aquele em que todos os testes continuam passando mesmo com o código alterado; nesse caso, a suíte não conseguiu distinguir o mutante do programa original, o que sugere lacuna na cobertura ou na assertividade dos testes. O objetivo é ter uma suíte que mate o máximo possível de mutantes (alto Mutation Score).

Foi criada uma cópia do código em \texttt{caixa\_eletronico\_mutante.py}
com mutantes manuais (ex.: permitir depósito zero; impedir saque de todo o saldo;
transferência na direção errada; extrato em ordem invertida).

Fórmula do Mutation Score:
\[
MS = \frac{\text{mutantes mortos}}{\text{mutantes totais} - \text{mutantes equivalentes}} \times 100\%.
\]

\begin{table}[H]
  \centering
  \caption{Resultados dos mutantes manuais}
  \label{tab:mutantes-manuais}
  \begin{tabular}{lcc}
    \toprule
    & Quantidade & MS (\%) \\
    \midrule
    Mutantes totais                & (preencher) & \\
    Mutantes equivalentes           & (preencher) & \\
    Mutantes mortos                 & (preencher) & \\
    Mutantes sobreviventes          & (preencher) & \\
    \midrule
    \textbf{Mutation Score (MS)}    & & (preencher) \\
    \bottomrule
  \end{tabular}
\end{table}

\section{Mutação Automática com \texttt{mutmut}}

Foi utilizada a ferramenta \texttt{mutmut} (configurada em \texttt{pyproject.toml}).
Comandos: \texttt{mutmut run} e \texttt{mutmut results}.

\begin{table}[H]
  \centering
  \caption{Resultados do \texttt{mutmut}}
  \label{tab:mutmut}
  \begin{tabular}{lcc}
    \toprule
    & Quantidade & MS (\%) \\
    \midrule
    Mutantes totais                & (preencher) & \\
    Mutantes equivalentes           & (preencher) & \\
    Mutantes mortos                 & (preencher) & \\
    Mutantes sobreviventes          & (preencher) & \\
    \midrule
    \textbf{Mutation Score (MS)}    & & (preencher) \\
    \bottomrule
  \end{tabular}
\end{table}

\section{Discussão e Conclusão}

(Comparar resultados dos mutantes manuais com os do \texttt{mutmut};
discutir mutantes sobreviventes e possíveis novos testes.)

\end{document}
